\documentclass[8pt,a4paper]{extarticle}
\usepackage[T1]{fontenc}
\usepackage{lmodern}
\usepackage[a4paper,top=12mm,bottom=7mm,left=8mm,right=8mm,
            headheight=10pt,headsep=3mm,includehead]{geometry}
\usepackage{multicol,listings,xcolor,fancyhdr,amsmath,amssymb}
\usepackage{titlesec,enumitem,booktabs}
\usepackage[hidelinks]{hyperref}

% ---- ICPC TRD rules: <= 25 pages, single-sided, A4,
% ---- page number top-right, university + team top-left.
\pagestyle{fancy}\fancyhf{}
\fancyhead[L]{\small\bfseries Jagiellonian University --- team UZI}
\fancyhead[C]{\small\nouppercase{\leftmark}}
\fancyhead[R]{\small\bfseries \thepage}
\renewcommand{\headrulewidth}{0.4pt}
\renewcommand{\sectionmark}[1]{\markboth{#1}{}}

\setlength{\columnsep}{5mm}
\setlength{\columnseprule}{0.3pt}
\setlength{\parindent}{0pt}
\setlength{\parskip}{1pt}
\titleformat{\section}{\large\bfseries}{\thesection}{0.5em}{}
  [{\titlerule[0.6pt]}]
\titlespacing*{\section}{0pt}{4pt}{2pt}
\titleformat{\subsection}{\normalsize\bfseries}{}{0em}{}
\titlespacing*{\subsection}{0pt}{4pt}{1pt}
\setlist{nosep,leftmargin=1.2em}

\lstset{
  language=C++, basicstyle=\fontsize{8}{9}\selectfont\ttfamily,
  columns=fixed, basewidth=4.2pt, keepspaces=true,
  keywordstyle=\bfseries, commentstyle=\itshape\color{black!70},
  showstringspaces=false, tabsize=2, breaklines=true,
  breakatwhitespace=false, aboveskip=1pt, belowskip=2pt,
  frame=none, xleftmargin=0pt, upquote=true,
  literate={~}{{\textasciitilde}}1,
}
% \snipn{note}{Title}{file}: note = when / how / pitfalls
\newcommand{\snip}[2]{\subsection[#1]{#1 \hfill
  {\normalfont\small\ttfamily\detokenize{#2}}}%
  \lstinputlisting{src/#2}}
\newcommand{\snipn}[3]{\subsection[#2]{#2 \hfill
  {\normalfont\small\ttfamily\detokenize{#3}}}%
  {\small\setlength{\parskip}{0pt}#1\par}%
  \lstinputlisting{src/#3}}
\newcommand{\W}{\textbf{When:} }
\newcommand{\Ho}{\textbf{How:} }
\newcommand{\Pf}{\textbf{Pitfalls:} }
\newcommand{\shsnip}[2]{\subsection{#1}%
  \lstinputlisting[language=bash]{src/#2}}

\begin{document}
\begin{multicols*}{2}
{\Large\bfseries ICPC Team Reference Document}\\[1pt]
{\small Jagiellonian University, Krak\'ow --- team UZI}
\vspace{-6pt}
{\small\renewcommand{\contentsname}{}
\setcounter{tocdepth}{2}
\makeatletter\renewcommand{\@dotsep}{1000}\makeatother
\makeatletter\renewcommand{\@tocrmarg}{1.6em}\renewcommand{\@pnumwidth}{1em}\makeatother
\begin{multicols}{2}\raggedright\tableofcontents\end{multicols}}

\section{Contest}
\snip{Template}{0-contest/template.cpp}
\shsnip{Compile \& stress test}{0-contest/commands.sh}
\snip{Fast input}{0-contest/fastio.cpp}

\subsection{Limits ($\approx 10^8$ simple operations/s)}
{\small
\begin{tabular}{@{}ll|ll@{}}
$n \le 10$--$11$ & $n!$, $n!\cdot n$ & $n \le 10^5$ & $n\sqrt n$\\
$n \le 20$ & $2^n n$ & $n \le 10^6$ & $n \log n$\\
$n \le 500$ & $n^3$ & $n \le 10^8$ & $n$ (simple)\\
$n \le 5000$ & $n^2$ & $n \le 10^{18}$ & $\log n$, $\sqrt[3]{n}$\\
\end{tabular}\par
\texttt{int} $< 2.1\cdot10^9$, \texttt{ll} $< 9.2\cdot10^{18}$,
\texttt{\_\_int128} $< 1.7\cdot 10^{38}$. \texttt{double}: 15 digits;
\texttt{long double}: 18 digits.
}

\subsection{Before submitting / after WA}
\begin{itemize}
\item Re-read the statement: output format, 0/1-indexing, ``modulo'',
  multiple test cases (reset all global state!), limits of \emph{every} input.
\item Overflow: \texttt{int*int}, sums, \texttt{1<<40} (use \texttt{1LL<<}),
  \texttt{INF} too small / \texttt{INF+INF}.
\item Edge cases: $n=1$, $n=0$, all equal, already sorted, negative numbers,
  maximum test (time!), disconnected graph, self-loops, multi-edges.
\item Uninitialised variables, \texttt{<} vs \texttt{<=}, reading the whole
  input in multitest even when the answer is known early.
\item Recursion depth $10^5$--$10^6$ may crash on an 8 MB stack. Locally: \texttt{ulimit -s unlimited}. Judges (DOMjudge) usually allow stack = memory limit; if unsure, go iterative / BFS order.
\item Precision: compare with \texttt{EPS}, print \texttt{setprecision(10)};
  avoid \texttt{sqrt}/division in integer geometry.
\item TLE: \texttt{endl} in loops (use \texttt{'\textbackslash n'}),
  \texttt{map}$\to$sort+vector, pass vectors by reference,
  \texttt{memset} per test case, \texttt{\#pragma GCC optimize("O3")}.
\item Write a brute force + generator and stress test (above).
  Small cases also reveal formulas: print answers for $n=1..10$ and look
  for a pattern.
\end{itemize}

\subsection{When stuck: ideas to try}
\begin{itemize}
\item Brute force small cases and look at the numbers (pattern, OEIS-like
  formula, monotonicity $\Rightarrow$ binary search on the answer).
\item Sort + greedy; prove by exchange argument. Process in reverse
  (deletions $\to$ insertions with DSU). Fix one element / the maximum /
  the answer value and count the rest.
\item Contribution technique: $\sum$ over pairs/subarrays $=\sum$ over
  elements of (value $\times$ number of structures where it matters).
\item Complement counting; inclusion--exclusion over ``bad'' events.
\item Invariants and parity (sum mod $k$, number of inversions, colouring
  of a grid). Constructive: find a pattern for small $n$, then induct.
\item Model as a graph: states = vertices (BFS/Dijkstra on (position,
  extra state)); constraints = flow / matching / 2-SAT.
\item Small parameter: bitmask over it, $\sqrt{n}$ split (heavy/light),
  meet in the middle, randomisation (random shuffle, random pivot,
  hashing, ``answer is one of few candidates'').
\item Sweep line + data structure over events sorted by coordinate.
\item Binary search on the answer + greedy/flow check; parallel binary
  search for many queries.
\end{itemize}

\section{Mathematics}
\snipn{\W any answer ``modulo $p$''. Change \texttt{MOD} for $10^9+7$; division needs a prime \texttt{MOD}. \Pf \texttt{mint(x)} reduces any \texttt{ll} (negatives OK); compare with \texttt{==}, never with \texttt{<}.}{Modular integer (998244353)}{1-math/mint.cpp}
\snipn{\Ho call \texttt{initFact()} once; \texttt{N} must exceed the largest $n$ used (e.g.\ $n+k$ in stars and bars).}{Binomial coefficients}{1-math/binom.cpp}
\snipn{\W modulus not prime, linear congruences. $ax\equiv b \pmod m$: solvable iff $g=\gcd(a,m)\mid b$; then $x = (b/g)\,x_0 \bmod (m/g)$ with $x_0$ from \texttt{egcd(a, m, x0, y0)}. CRT for many congruences: fold pairwise.}{Power, extended gcd, inverse, CRT}{1-math/modpow.cpp}
\snip{Sieve, factorization, divisors, phi}{1-math/sieve.cpp}
\snip{Miller--Rabin (64-bit)}{1-math/millerrabin.cpp}
\snipn{\W factorise numbers up to $10^{18}$ (trial division would need $10^9$ steps). \Ho rho finds a non-trivial divisor, recurse on both parts. Result unsorted: \texttt{sort} it, then group equal primes.}{Pollard rho}{1-math/pollard.cpp}
\snipn{\W linear recurrence / DP with fixed transitions and huge $n$ (state size $k\le 100$); walks of length $n$ = $A^n$. \Ho state vector $v_{n+1} = M v_n$, answer $M^n v_0$. Prefix sums: add an accumulator state.}{Matrix power}{1-math/matrix.cpp}
\snipn{\W systems of linear equations, rank, counting solutions ($p^{m-\mathrm{rank}}$). Over GF(2): rows as \texttt{bitset}, subtraction is XOR (same code, 64$\times$ faster). Real numbers: pick the row with the largest $|a_{ic}|$ as pivot, compare with EPS.}{Gaussian elimination mod $p$}{1-math/gauss.cpp}
\snip{Determinant mod $p$}{1-math/det.cpp}
\snipn{\W max / count / existence of subset XOR. Two bases merge by inserting all vectors of one into the other ($O(60^2)$).}{XOR basis}{1-math/xorbasis.cpp}
\snipn{\W polynomial products, ``number of pairs with $a_i+b_j=s$'' for all $s$ (convolve count arrays), DP transitions that are convolutions. \Pf inputs must be in $[0,\mathrm{MOD})$; other modulus or exact big counts $\Rightarrow$ FFT below.}{NTT (convolution mod 998244353)}{1-math/ntt.cpp}
\snipn{\W exact integer convolution without modulus (results up to $10^{15}$), e.g.\ counting pairs by sum, string matching with wildcards ($\sum (a-b)^2ab$ trick).}{FFT (exact integer convolution)}{1-math/fft.cpp}
\snipn{\W the answer for $n$ satisfies an (unknown) linear recurrence: counting with a fixed-width automaton/DP, tilings, $a_n = \sum c_j a_{n-j}$. \Ho brute-force the first $2L+10$ terms (e.g.\ 500), run \texttt{berlekampMassey}, then \texttt{linearRec} for $n$ up to $10^{18}$. \Pf works only modulo a prime; give it enough terms.}{Berlekamp--Massey, $k$-th term}{1-math/bm.cpp}
\snipn{\W answer is a polynomial in $n$ of known degree $d$ (sums of powers, counting in $d$ nested loops): compute $y_0..y_d$ by brute force, evaluate at $n$.}{Lagrange interpolation}{1-math/lagrange.cpp}
\snipn{\W $a^x\equiv b$; floor sums appear in counting lattice points under a line and in $\sum \lfloor (ai+b)/m\rfloor$-type formulas.}{Discrete log, floor sum}{1-math/dlog.cpp}

\subsection{Number theory facts}
\begin{itemize}
\item Primes: $998244353 = 119\cdot2^{23}+1$ (root 3),
  $10^9+7$, $10^9+9$, $10^{18}+3$; $2^{61}-1$.
\item Fermat: $a^{p-1}\equiv1 \pmod p$; Euler: $a^{\varphi(m)}\equiv 1$
  if $\gcd(a,m)=1$. For huge exponents: $a^e \equiv a^{e \bmod (p-1)} \pmod p$
  ($p\nmid a$).
\item Max number of divisors: $\le 10^6$: 240; $\le10^9$: 1344;
  $\le 10^{18}$: 103680. $\pi(10^6)=78498$, $\pi(10^9)\approx 5.08\cdot10^7$.
  Prime gaps $< 300$ below $10^{18}$.
\item $\sum_{d\mid n}\varphi(d)=n$. $\gcd(F_a,F_b)=F_{\gcd(a,b)}$.
  $\gcd(a,b)=\gcd(a,b-a)$.
  $\sum_{i\le n}\lfloor n/i\rfloor$ has $O(\sqrt n)$ distinct values:
  \texttt{for (l=1; l<=n; l=r+1) \{ r = n/(n/l); ... \}}.
\item Möbius: $\sum_{d\mid n}\mu(d)=[n=1]$; number of coprime pairs
  $(a,b)\le n$: $\sum_d \mu(d)\lfloor n/d\rfloor^2$.
\item Sum of divisors $\sigma(n)=\prod \frac{p^{e+1}-1}{p-1}$;
  number of divisors $\prod (e+1)$.
\item Wilson: $(p-1)!\equiv -1 \pmod p$. Lucas: $\binom{n}{k}\equiv\prod
  \binom{n_i}{k_i} \pmod p$ (base-$p$ digits).
  $\binom{n}{k}$ is odd $\iff k \mathbin{\&} n = k$.
\end{itemize}

\subsection{Combinatorics}
\begin{itemize}
\item $\binom{n}{k}=\binom{n-1}{k-1}+\binom{n-1}{k}$;
  $\sum_{i=0}^{n}\binom{i}{k}=\binom{n+1}{k+1}$;
  $\sum_k\binom{m}{k}\binom{n}{r-k}=\binom{m+n}{r}$ (Vandermonde).
\item Stars and bars: $x_1+\dots+x_k=n$, $x_i\ge0$:
  $\binom{n+k-1}{k-1}$; with $x_i\ge1$: $\binom{n-1}{k-1}$.
\item Catalan $C_n=\frac{1}{n+1}\binom{2n}{n}=\binom{2n}{n}-\binom{2n}{n+1}$:
  1, 1, 2, 5, 14, 42, 132, 429, 1430 (bracket sequences, binary trees,
  triangulations, monotone paths not crossing the diagonal).
  Paths $(0,0)\to(a,b)$ touching line $y=x+1$: reflect $\Rightarrow
  \binom{a+b}{a+1}$.
\item Derangements $D_n=(n-1)(D_{n-1}+D_{n-2})$, $D_n\approx n!/e$: 1, 0, 1, 2, 9, 44, 265.
\item Stirling 2nd kind $S(n,k)=kS(n-1,k)+S(n-1,k-1)$ (partitions into
  $k$ non-empty sets); surjections $=k!\,S(n,k)=\sum_i(-1)^i\binom{k}{i}(k-i)^n$.
  Bell: 1, 1, 2, 5, 15, 52, 203, 877.
\item Stirling 1st kind (permutations with $k$ cycles):
  $c(n,k)=(n-1)c(n-1,k)+c(n-1,k-1)$.
\item Integer partitions $p(n)$: 1, 1, 2, 3, 5, 7, 11, 15, 22, 30, 42;
  $p(100)\approx 1.9\cdot10^8$.
\item Inclusion--exclusion: $|\bigcup A_i|=\sum_{\emptyset\neq S}(-1)^{|S|+1}|\bigcap_{i\in S}A_i|$.
\item Burnside: \#orbits $=\frac{1}{|G|}\sum_{g}\mathrm{fix}(g)$; necklaces
  of $n$ beads, $k$ colours (rotations): $\frac1n\sum_{d\mid n}\varphi(d)k^{n/d}$.
\item Cayley: $n^{n-2}$ labelled trees; forests of $k$ trees rooted at
  given $k$ vertices: $k n^{n-k-1}$. Kirchhoff: \#spanning trees $=$ any
  cofactor of $D-A$.
\item Expected value is linear: $E[\sum X_i]=\sum E[X_i]$ (indicators!).
  Geometric: success prob.\ $p$ $\Rightarrow$ expected $1/p$ tries.
  Coupon collector: $n H_n$. $P(A)$ as fraction: $a\cdot b^{-1} \bmod p$.
\end{itemize}

\subsection{Sums}
$\sum i=\frac{n(n+1)}2$, $\sum i^2=\frac{n(n+1)(2n+1)}6$,
$\sum i^3=\left(\frac{n(n+1)}2\right)^2$,
$\sum_{i=0}^{n} r^i=\frac{r^{n+1}-1}{r-1}$, $H_n\approx\ln n+0.5772$,
$\sum_{p\le n} 1/p \approx \ln\ln n$.

\subsection{Games}
Sprague--Grundy: $g(v)=\mathrm{mex}\{g(u)\}$ over moves;
a sum of games is lost $\iff$ XOR of $g=0$. Nim: XOR of piles.
Staircase Nim: XOR of odd steps. Misère Nim: as normal unless all piles
$\le1$ (then win $\iff$ even number of piles of 1).
Typical trick: compute $g$ by brute force for small states, find the pattern.

\section{Data structures}
\snip{Disjoint set union}{2-ds/dsu.cpp}
\snip{Fenwick tree}{2-ds/fenwick.cpp}
\snipn{\Ho change \texttt{T}, \texttt{ID} and \texttt{f}: sum ($0$), min (\texttt{LLONG\_MAX}), max (\texttt{LLONG\_MIN}), gcd ($0$), or a struct (e.g.\ max subarray sum: \{sum, pref, suf, best\}). \texttt{f} must be associative, need not be commutative. Point add: \texttt{set(i, t[i + n] + v)}.}{Segment tree (point update)}{2-ds/segtree.cpp}
\snipn{\Ho a node stores the answer for its range plus a pending tag. \texttt{apply} updates both; \texttt{push} hands the tag to the children before going down. Other operations: change the tag. Range assign: tag = (has, value); assign overrides older tags. Assign + add: compose tags in order (new assign clears add).}{Lazy segment tree (range add)}{2-ds/lazyseg.cpp}
\snipn{\W static array, many min/max/gcd queries (idempotent operations only; for sums use prefix sums).}{Sparse table}{2-ds/sparse.cpp}
\snip{Monotonic stack}{2-ds/monostack.cpp}
\snip{Ordered set, hash map}{2-ds/pbds.cpp}
\snipn{\W DP of the form $dp_i = \min_j (k_j x_i + b_j)$ (convex hull trick) with lines and queries in any order. \Ho add the line of $j$ after computing $dp_j$; query at $x_i$. Range of $x$ is the integer domain $[lo, hi)$; nodes are created on demand, so $\pm 10^9$ is fine. \Pf $k\cdot x+b$ must fit in \texttt{ll}.}{Li Chao tree (lines, min)}{2-ds/lichao.cpp}
\snipn{\W an array with insert / erase / reverse / cut-and-paste at arbitrary positions plus range aggregates. \Ho every operation = \texttt{split} into pieces, modify a piece (set a tag), \texttt{merge} back. More tags: handle them in \texttt{push} like \texttt{Rev}, recompute in \texttt{pull}.}{Implicit treap}{2-ds/treap.cpp}
\snipn{\W offline dynamic connectivity (edges added and removed): each edge lives on a time interval; put it into the $O(\log q)$ segment-tree nodes over time covering it; DFS that tree: unite the node's edges, recurse, \texttt{rollback}. Also bipartiteness over time (DSU with parity).}{DSU with rollback}{2-ds/dsu_rollback.cpp}
\snipn{\W static array, many range queries whose answer is easy to update when one element enters/leaves the window (number of distinct values, pairs of equal values, mode with counts). Offline only. $n=q=2\cdot10^5$: about $10^8$ add/rem calls.}{Mo's algorithm}{2-ds/mo.cpp}
\snipn{\W queries on old versions or on ``prefix $r$ minus prefix $l$'': $k$-th smallest in a range, count of values $\le x$ in a range (sum \texttt{Cnt} of left children while descending).}{Persistent segment tree}{2-ds/persistent.cpp}
\subsection{Ideas}
\begin{itemize}
\item Offline: sort queries by right end + Fenwick/segment tree
  (``number of distinct values on $[l,r]$'': keep only the last
  occurrence). Process in reverse to turn deletions into DSU unions.
\item Sqrt decomposition: blocks of $\approx\sqrt n$; Mo's order for
  $[l,r)$ queries: sort by $(l/B,\ r)$, alternate $r$ direction for odd
  blocks; $B\approx n/\sqrt q$.
\item Small-to-large merging of sets: $O(n\log^2 n)$ total.
\item Segment tree descent: first index with prefix $\ge x$ in $O(\log n)$
  (walk down from the root).
\item Persistent: copy the $O(\log n)$ nodes on the path on each update.
\end{itemize}

\section{Graphs}
\snip{Dijkstra, 0-1 BFS, Bellman--Ford}{3-graph/shortest.cpp}
\snip{Topological sort}{3-graph/topo.cpp}
\snipn{\Ho 2-SAT: write every constraint as clauses ``$a$ or $b$''. ``$a \Rightarrow b$'' $=$ \texttt{either(a false, b true)}; ``$a = b$'' = two clauses; ``exactly one'' = $(a\lor b)\land(\lnot a\lor\lnot b)$. SCC condensation is a DAG: DP on it in order $ncomp-1 \to 0$ (comp 0 is a sink).}{SCC (Tarjan) and 2-SAT}{3-graph/scc.cpp}
\snipn{\W critical edges/vertices. Contract 2-edge-connected components (DFS without bridges) $\Rightarrow$ the \emph{bridge tree}; many problems become tree problems on it.}{Bridges and articulation points}{3-graph/bridges.cpp}
\snipn{\W distances/path sums on a static tree: $f(u,v) = p_u + p_v - 2p_{lca}$ for root-prefix values $p$. \texttt{kth} gives the $k$-th vertex on a path.}{LCA (binary lifting)}{3-graph/lca.cpp}
\snipn{\W path update + path query (sum/max/assign) on a tree that changes values: $O(\log^2 n)$ per query with a segment tree on \texttt{pos}. Subtree query = one range. \texttt{pos} is a DFS preorder. \Pf recursion depth $n$.}{Heavy-light decomposition}{3-graph/hld.cpp}
\snipn{\W counting/optimising over \emph{all paths} (number of paths of length $\le K$, with sum $=X$), or ``nearest marked vertex'' queries. \Ho every path is handled at the highest centroid on it; at centroid $c$, DFS each child subtree, combine with data from previous subtrees (avoid pairs inside one subtree). Nearest marked: \texttt{best[c]} for every centroid ancestor $c$ of $v$; answer $=\min_c$ \texttt{best[c]} $+\,dist(v,c)$; depth $\le\log n$.}{Centroid decomposition}{3-graph/centroid.cpp}
\snipn{\W many queries, each on a small set of $k$ vertices with $\sum k$ bounded: build a tree of size $<2k$ and run the DP only there.}{Virtual tree}{3-graph/virtualtree.cpp}
\snipn{\W ``for every subtree, something about the multiset of values in it''. Same idea for merging heaps, maps of counts, etc.}{Small-to-large (DSU on tree)}{3-graph/sack.cpp}
\snip{Euler path}{3-graph/euler.cpp}
\snipn{\W bipartite matching up to $\sim 5000$ vertices per side; larger: Dinic ($O(E\sqrt V)$). Shuffle the adjacency lists against bad cases.}{Bipartite matching (Kuhn)}{3-graph/matching.cpp}
\snipn{\Ho modelling: source $\to$ suppliers, consumers $\to$ sink, capacities = limits. Undirected edge: \texttt{addEdge(u, v, c, c)}. Vertex capacity: split $v$ into $v_{in}\to v_{out}$. Min cut edges = saturated edges from \texttt{lvl}$\ge 0$ to \texttt{lvl}$<0$. \Pf \texttt{s != t}.}{Max flow (Dinic)}{3-graph/dinic.cpp}
\snipn{\W flow with costs: assignment, transport, ``choose $k$ disjoint paths of min total length''. \Ho augments along cheapest paths, so with \texttt{lim = K} you get the cheapest $K$ units. \Pf negative cycles are not allowed (negative edges are).}{Min-cost flow}{3-graph/mcmf.cpp}
\snipn{\W dense assignment problem, $n\le 500$ ($n^2m$). Simpler and faster than min-cost flow for this case.}{Hungarian algorithm}{3-graph/hungarian.cpp}
\subsection{Facts}
\begin{itemize}
\item Tree: $n-1$ edges, connected. Diameter: farthest vertex $a$ from any
  vertex, then farthest $b$ from $a$. Center(s) on the diameter middle.
\item Euler: undirected connected graph has an Euler cycle $\iff$ all degrees
  even; path $\iff$ 0 or 2 odd vertices. Directed: in $=$ out everywhere
  (path: one vertex out$-$in$=1$, one in$-$out$=1$).
\item Planar: $V-E+F=2$ (connected), $E\le3V-6$. Cactus: \#cycles $=E-V+1$.
\item Bipartite $\iff$ no odd cycle. König: max matching $=$ min vertex cover.
  Hall: perfect matching of $L$ $\iff$ $|N(S)|\ge|S|$ for all $S\subseteq L$.
  Min path cover in a DAG $=n-$ max matching (split $v$ into $v_{out},v_{in}$).
  Dilworth: min chains $=$ max antichain.
\item Max flow $=$ min cut. Vertex capacities: split $v\to v_{in},v_{out}$.
  Lower bounds $l$: add demand edges $S'\to v$, $u\to T'$.
  Max closure (project selection): $\sum$ positive $-$ mincut
  (profit: $s\to v$, cost: $v\to t$, dependency $u\to v$ with $\infty$).
\item MST: Kruskal (sort + DSU) or Prim $O(m\log n)$; minimum
  bottleneck path lies on the MST. Boruvka: $O(\log n)$ rounds, each vertex
  picks its cheapest outgoing edge.
\item Shortest paths with $k$ free edges / states: graph on (vertex, state).
  Dijkstra with potentials for negative edges after one Bellman--Ford.
\item Functional graph ($n$ vertices, out-degree 1): every component is
  a cycle with trees; use binary lifting for $k$-th successor.
\item Tournament: has a Hamiltonian path; SCCs are linearly ordered.
\end{itemize}

\section{Strings}
\snip{Prefix function, Z-function}{4-string/kmp_z.cpp}
\snipn{\W compare substrings in $O(1)$: LCP by binary search, palindrome check (hash of reversed string), counting distinct substrings of one length. Collision probability $\approx n^2/2^{61}$: safe.}{Polynomial hashing}{4-string/hash.cpp}
\snip{Manacher}{4-string/manacher.cpp}
\snipn{\W sorted suffixes: find pattern by binary search ($O(|t|\log n)$), longest repeated substring $=\max$ lcp, LCS of $s,t$: SA of $s\#t$, max lcp of adjacent suffixes from different strings, $k$-th distinct substring.}{Suffix array + LCP}{4-string/sufarr.cpp}
\snipn{\W many patterns in one text; DP over automaton states (count strings of length $n$ avoiding all patterns: states $\times$ length, forbid states with \texttt{out > 0}).}{Aho--Corasick}{4-string/aho.cpp}
\snip{Binary trie (max XOR)}{4-string/trie.cpp}
\snipn{\W everything about all substrings in $O(n)$: distinct count, occurrences of each, LCS of two strings (walk $t$ on SAM($s$), on mismatch follow \texttt{link}, keep current length), smallest cyclic shift (SAM of $s+s$).}{Suffix automaton}{4-string/sam.cpp}
\subsection{Ideas}
\begin{itemize}
\item LCP of suffixes $i,j$ $=$ min of \texttt{lcp} on
  $[\mathrm{rk}_i,\mathrm{rk}_j)$ (sparse table).
\item Minimal rotation of $s$: suffix array / Booth on $s+s$.
  $s$ is a power of a word of length $n-p_{n-1}$ iff that divides $n$.
\item Count occurrences of every prefix: \texttt{cnt[p[i]]++}, then for $i$ from $n$
  down: \texttt{cnt[p[i-1]] += cnt[i]}, then \texttt{cnt[i]++}.
\end{itemize}

\section{Geometry}
\snip{Point}{5-geom/point.cpp}
\snip{Segments and lines}{5-geom/segments.cpp}
\snip{Polygons, convex hull}{5-geom/polygon.cpp}
\snip{Polar angle sort}{5-geom/angle.cpp}
\snip{Circles}{5-geom/circle.cpp}
\snip{Minimum enclosing circle}{5-geom/mec.cpp}
\snip{Point in convex polygon, closest pair}{5-geom/convex.cpp}
\snip{Polygon inside polygon}{5-geom/polyinside.cpp}
\subsection{Formulas}
\begin{itemize}
\item Triangle: area $=\frac12|ab\sin C|=\sqrt{s(s-a)(s-b)(s-c)}$;
  $R=\frac{abc}{4S}$, $r=\frac{S}{s}$; law of cosines
  $c^2=a^2+b^2-2ab\cos C$; $\frac{a}{\sin A}=2R$.
\item Line through $p,q$: $a=q_y-p_y$, $b=p_x-q_x$, $c=-(ap_x+bp_y)$;
  distance from $r$: $|ar_x+br_y+c|/\sqrt{a^2+b^2}$.
\item Rotation by $\theta$: $(x\cos\theta-y\sin\theta,\ x\sin\theta+y\cos\theta)$.
\item Circle segment area: $\frac{r^2}{2}(\theta-\sin\theta)$.
  Sphere: $V=\frac43\pi r^3$, $S=4\pi r^2$; cone $V=\frac13\pi r^2h$.
\item Pick: $A=I+\frac B2-1$. Lattice points on segment: $\gcd(|dx|,|dy|)+1$.
\item Manhattan distance: rotate $(x+y,\,x-y)$, then Chebyshev
  ($\max$ of coordinate differences).
\item Closest pair: sort by $x$, sweep keeping a set ordered by $y$
  within the current best distance $d$; $O(n\log n)$.
\item Points in convex polygon in $O(\log n)$: binary search on the fan
  from vertex 0.
\end{itemize}

\section{DP and techniques}
\snipn{\W D\&C optimisation: $dp_{new}[i] = \min_{j<i} dp_{old}[j] + C(j,i)$ and the best $j$ is non-decreasing in $i$ (true if $C(a,c)+C(b,d)\le C(a,d)+C(b,c)$ for $a\le b\le c\le d$, e.g.\ $C$ = cost of a segment that grows superadditively). Splitting into $k$ groups: $O(kn\log n)$.}{DP patterns}{6-misc/dp.cpp}
\snip{Knuth optimisation}{6-misc/knuth.cpp}
\snip{Aliens trick (exactly $K$ items)}{6-misc/aliens.cpp}
\snip{Re-rooting}{6-misc/reroot.cpp}
\snip{Binary / ternary search}{6-misc/search.cpp}
\snip{Bits, random, misc}{6-misc/tricks.cpp}
\snipn{\W area/volume with curved boundaries (union of circles via $\int$ length of the vertical cut). Split at non-smooth points.}{Numerical integration}{6-misc/simpson.cpp}
\subsection{DP optimisations}
\begin{itemize}
\item Convex hull trick: $dp_i=\min_j(k_j x_i+b_j)$ $\Rightarrow$ Li Chao tree.
\item Knuth: $dp[l][r]=\min_{l\le m<r} dp[l][m]+dp[m+1][r]+w(l,r)$ with
  $opt[l][r-1]\le opt[l][r]\le opt[l+1][r]$ $\Rightarrow O(n^2)$.
\item ``Aliens'' / Lagrangian relaxation: exactly $k$ items with convex
  answer in $k$ $\Rightarrow$ binary search a penalty $\lambda$ per item.
\item Exchange argument for greedy order: compare $a$ before $b$ vs.\ $b$
  before $a$ and sort by that comparator (must be a strict weak order).
\item Tree DP: $dp[v]$ from children; rerooting: second pass from the root
  passing the ``parent side'' value down (prefix/suffix of children).
\item Small answers/strange constraints ($n\le 20$, $k\le 10$, $\sum\le 10^5$)
  hint at bitmask, meet-in-the-middle ($2^{n/2}$), or $\sqrt{}$ splitting.
\end{itemize}
\subsection{STL reminders and pitfalls}
\begin{itemize}
\item \texttt{s.lower\_bound(x)} on \texttt{set}/\texttt{map} is $O(\log n)$;
  \texttt{lower\_bound(all(s), x)} on a set is $O(n)$!
\item \texttt{accumulate(all(v), 0LL)} (with \texttt{0} it sums in \texttt{int}).
  \texttt{abs} on \texttt{ll}: \texttt{llabs} or \texttt{abs((ll)x)}.
\item \texttt{v.size() - 1} is unsigned: \texttt{sz(v) - 1}. \texttt{(int)sqrt(x)} may be off by one:
  adjust with \texttt{while (r * r > x) r--; while ((r+1)*(r+1) <= x) r++;}.
\item \texttt{nth\_element(v.begin(), v.begin() + k, v.end())}: $k$-th in $O(n)$.
  \texttt{partial\_sum}, \texttt{iota}, \texttt{\_\_gcd}/\texttt{gcd}, \texttt{lcm} (overflow!).
\item \texttt{multiset::erase(x)} erases ALL copies; use \texttt{s.erase(s.find(x))}.
\item \texttt{priority\_queue} is a max-heap; min-heap: \texttt{greater<>}.
  Sorting with a comparator: must be a strict weak order (\texttt{<}, never \texttt{<=}).
\item \texttt{map[x]} inserts $x$; use \texttt{count}/\texttt{find} to test.
  \texttt{unordered\_map}: \texttt{reserve} and \texttt{SafeHash}.
\item 2D prefix sums: $S_{i+1,j+1} = a_{ij} + S_{i,j+1} + S_{i+1,j} - S_{ij}$;
  difference array for range add: $d_l \mathrel{+}= v$, $d_r \mathrel{-}= v$, then prefix sum.
\item Interactive: flush after every query (\texttt{endl} or \texttt{cout.flush()});
  do not use \texttt{fread} input; stop on \texttt{-1}.
\end{itemize}
\end{multicols*}
\end{document}
